\documentclass[10pt,a4paper]{amsart}
\usepackage[margin=19mm]{geometry}
\usepackage{amsmath,amssymb,mathtools}
\usepackage[scr=rsfs,cal=euler]{mathalfa}
\usepackage{xfrac}
\usepackage[unicode,hidelinks,bookmarks=false]{hyperref}
\newcommand{\ie}{i.e.\ }
\newcommand{\cf}{cf.\ }
\newcommand{\resp}{respectively\ }
\newcommand{\Subsection}{\subsection}
\providecommand{\nobreakdash}{\nobreak\hbox{-}}
\setlength{\emergencystretch}{2em}
\multlinegap0pt
\usepackage{float}
\usepackage{multirow}
\usepackage{amssymb}
\usepackage{xcolor}
\usepackage{booktabs}
\usepackage{algorithm}\usepackage{algpseudocode}
\newtheorem{lemma}{Lemma}
\title{On the supremum and its location of the standardized uniform empirical process}
\author{Dietmar Ferger}
\date{Source: arXiv:2512.17674v1 (CC BY 4.0)}
\begin{document}
\maketitle

\noindent\emph{Source and licence.} On the supremum and its location of the standardized uniform empirical process. Version arXiv:2512.17674v1 (CC BY 4.0), distributed by arXiv under the Creative Commons Attribution 4.0 International licence, http://creativecommons.org/licenses/by/4.0/, as recorded in the arXiv licence metadata for this exact version and shown on the abstract page https://arxiv.org/abs/2512.17674v1. Full licence notice: \texttt{licenses/arxiv-2512.17674-CC-BY-4.0.txt}.\par\smallskip

\noindent\textit{Editorial scope.} Counted unit: the article's lemma maxinequalityQn (source0.tex lines 263--267). The article's complete original proof of it is delivered with it (source0.tex lines 269--283, one \texttt{proof} environment of 15 lines). No mathematics is written, repaired or re-derived: the statement and the proof are the article's own lines, in the article's own notation, and no retained line is substituted or re-flowed.  No further source-local context is needed: every cross-reference of every retained line resolves inside the two delivered spans.  Nothing was omitted from any retained span, so the package carries no omission record.  The retained lines cite 2 works, whose entries are transcribed verbatim from the article's own bibliography source in the e-print and delivered with short editorial labels recorded in the item's reference mapping.  No retained line contains a figure environment, an image-inclusion command or drawing code, so no image is delivered and the package declares no asset dependency.  Editorial formatting only, in addition: an AMS \texttt{amsart} preamble that restates the article's own declarations and macros used by the retained lines, namely the environments \texttt{lemma} on separate counters, together with the standard packages those lines need.  The retained environments print in this document as Lemma 1 in order of appearance, whereas the article prints the same items with the numbering of its own full document.  The complete native source of the article as distributed in the arXiv e-print for this exact version is bundled unchanged as \texttt{sources/191336/source0.tex}.
\begin{lemma} \label{maxinequalityQn} For every $a \in (0,\frac{1}{2}]$ and each $\lambda>0$ it follows that
\begin{equation} \label{maxinequality}
\mathbb{P}\big(\sup_{a \le t \le 1-a} \frac{|F_n(t)-t|}{\sqrt{t(1-t)}} > \lambda \big) \le 2 \; \lambda^{-2} \; n^{-1} \big(\log(\frac{1-a}{a})+1\big).
\end{equation}
\end{lemma}

\begin{proof} It is well-known that the process $\frac{F_n(t)-t}{1-t}, t \in [0,1),$ is a centered martingale with respect to the natural filtration, confer Gaenssler and Stute \cite{Gaenssler}, p.4. Therefore, $S(t):= (\frac{F_n(t)-t}{1-t})^2, t \in [0,1),$ is a non-negative sub-martingale.
From Lemma 3.3 in Ferger \cite{Ferger1} we can infer with $w(t)=\frac{1-t}{t}$ that

\begin{eqnarray}
\mathbb{P}\big(\sup_{a \le t \le 1/2} \frac{|F_n(t)-t|}{\sqrt{t(1-t)}} > \lambda \big)
&\le&\mathbb{P}\big(\sup_{a \le t \le 1/2} w(t) S(t) > \lambda^2 \big) \nonumber\\
&\le& \lambda^{-2}\{\int_a^{\frac{1}{2}} H(t) (-w)(dt)+w(\frac{1}{2}) H(\frac{1}{2})\}. \label{left}
\end{eqnarray}
Here, $H(t)=\mathbb{E}[S(t)]= (1-t)^{-2} \text{Var}[F_n(t)]= n^{-1} \frac{t}{1-t}.$ Since $-w^\prime(t)=t^{-2}$, the integral in (\ref{left}) simplifies to
$n^{-1} \log(\frac{1-a}{a})$ and $w(\frac{1}{2}) H(\frac{1}{2})=n^{-1}$. Thus we obtain from (\ref{left}) that
$$
 \mathbb{P}\big(\sup_{a \le t \le 1/2} \frac{|F_n(t)-t|}{\sqrt{t(1-t)}} > \lambda \big) \le \lambda^{-2} \; n^{-1} (\log(\frac{1-a}{a})+1).
$$
Now, the assertion follows, because $$\sup_{a \le t \le 1/2} \frac{|F_n(t)-t|}{\sqrt{t(1-t)}} \stackrel{\mathcal{D}}{=} \sup_{1/2 \le t \le 1-a} \frac{|F_n(t)-t|}{\sqrt{t(1-t)}}.$$
\end{proof}

\begin{thebibliography}{99}
\bibitem[Ferger]{Ferger1} D. Ferger, \textit{Density estimation via best $L^2$-approximation on classes of step functions}, Kybernetika \textbf{53} (2017), 198--219.
\bibitem[Gaenss]{Gaenssler} P. Gaenssler and W. Stute, \textit{Seminar on Empirical Processes}, Basel, Boston: Birkh\"{a}user, 1987.
\end{thebibliography}
\end{document}
